\section{Numerical method and validation}
\label{sec:numerics}

\subsection{Geometry and integration parametrization}
\label{sec:importance}

We study the seven geometries listed in \cref{tab:geometries}. The common
point $(5,1)$ links a scan of axial confinement to a scan of lateral
confinement.

\begin{table}[ht]
 \centering
 \caption{Quantum-dot semiaxes used in the axial and lateral scans, in
 effective Bohr radii.}
 \label{tab:geometries}
 \begin{tabular}{lcc}
  \toprule
  Scan & $a$ & $c$\\
  \midrule
  Axial & 5 & 0.5, 1, 1.5, 2\\
  Lateral & 3, 5, 7, 10 & 1\\
  \bottomrule
 \end{tabular}
\end{table}

Direct integration in the physical coordinates is inefficient because the
single-particle orbitals occupy only a small part of the nominal domain. We
therefore transform every particle to variables matched to the ground-state
density. The axial coordinate is first written
\begin{equation}
 z=u\,c\sqrt{1-r^2/a^2},\qquad 0<u<1 .
 \label{eq:axial-map}
\end{equation}
The numerical radial domain is restricted to the validated cutoff $R=0.9a$.
Using the ground-state in-plane length $\ell$, the reference radial and axial
densities are
\begin{equation}
 p_r(r)=\frac{1}{C}\frac{2r}{\ell^2}e^{-r^2/\ell^2},
 \quad C=1-e^{-R^2/\ell^2},
 \qquad
 p_u(u)=2\sin^2(\pi u).
 \label{eq:reference-densities}
\end{equation}
Both reference densities have analytical cumulative distributions. For the
radial density, the inverse is
\begin{equation}
 r(v)=\ell\sqrt{-\ln(1-Cv)},\qquad 0<v<1,
 \label{eq:radial-inverse}
\end{equation}
whereas the axial cumulative distribution has no elementary closed form inverse and is tabulated once and
interpolated. This transformation cancels most of the localized
single-particle density and leaves an order-unity transformed integrand. It is
an importance transformation, not a symmetry reduction.

The final estimator removes only the redundant global azimuthal rotation. For
the exciton, the hole angle is fixed to zero, the electron angle spans
$[0,2\pi]$, and a factor $2\pi$ restores the global rotation. For the
biexciton, $\phi_a=0$ while $\phi_1$, $\phi_2$, and $\phi_b$ are each
integrated independently over $[0,2\pi]$, again with an overall factor
$2\pi$. No reflection folding or particle-exchange multiplicities are used.
The complete $G(P+Q)$ factor, the six Coulomb interactions in
\cref{eq:biexciton-coulomb}, and the four kinetic contributions in
\cref{eq:biexciton-kinetic-numerator} remain in the pointwise integrand.

\subsection{Adaptive quasi-Monte Carlo quadrature}
\label{sec:quadrature}

Every integral is evaluated with the Wolfram Language implementation of
adaptive quasi-Monte Carlo quadrature. Low-discrepancy points improve the
coverage of the high-dimensional unit hypercube relative to independent
pseudorandom sampling \cite{caflisch1998}, while adaptive subdivision directs
additional evaluations toward regions with larger estimated error. The
production settings are
\begin{equation}
\begin{aligned}
 \texttt{Method}&=\texttt{AdaptiveQuasiMonteCarlo},\\
 \texttt{BisectionDithering}&=0,\\
 \texttt{AccuracyGoal}&=\infty,\qquad
 \texttt{PrecisionGoal}=4 .
\end{aligned}
\label{eq:qmc-settings}
\end{equation}
The stage-dependent point budget is supplied through
\texttt{MaxPoints}. A finite absolute accuracy goal is inappropriate because
some biexciton numerator and norm integrals are much smaller than a typical
absolute tolerance. Reaching \texttt{MaxPoints} therefore indicates that the
requested relative goal was not attained before the budget cap; the returned
integral value and internal error estimate are nevertheless retained.

The Hamiltonian numerator and norm denominator are independent adaptive calls.
They do not share a fixed Sobol point set. This preserves separate adaptation
to their different integrands. The internal quadrature errors returned for
the numerator and denominator are treated as independent and propagated to
the exciton and biexciton energies, the binding energy, and the radiative
lifetimes. The resulting error bars are numerical-resolution indicators, not
statistical confidence intervals: they neglect covariance between the
adaptive integrals, possible systematic quasi-Monte Carlo bias, and uncertainty
in the location of the variational optimum.

The final estimator was selected through a sequence of checks at $(5,1)$.
These included the exact exciton normalization at $\alpha_X=0$, integration
budget comparisons, dithered adaptive partitions, a complete four-angle
biexciton calculation, and the final global-rotation-fixed calculation. The
two angular formulations agree within the observed resolution, while removing
the redundant global angle reduces the numerical dimension. Earlier
direct/cross and particle-relabeling identities remain analytically valid,
but their independent numerical evaluations exhibited percent-level
differences and were therefore not used to assemble the final energies.

\subsection{Variational search}
\label{sec:optimization}

The variational search proceeded in two phases. The first was an exploratory
multidimensional search performed with the Wolfram Language function
\texttt{NMinimize} using the Nelder--Mead method. It varied the parameters
simultaneously at a lower integration budget and located approximate minima
for each geometry. Because small objective differences can be dominated by
quasi-Monte Carlo fluctuations, these first-phase optima were used only as
starting points.

The second phase refined those starting points with an automated local
coordinate pattern search using the integration setup of
\cref{sec:importance}. Positive and
negative displacements of each parameter were screened at a moderate budget;
the center and the most promising candidates were then recomputed at a larger
budget, with progressively smaller displacements used around stable
candidates. Automation generated and evaluated the candidates, but final
selection also relied on inspecting the energies, integration-error estimates,
stability between budgets, neighboring points in parameter space, and
consistency of the geometry trends. A parameter set was retained only when
this review indicated a stable local minimum rather than an isolated low
numerical fluctuation.
